\section{Discussion}
The final comparison clarifies the main narrative of the project more sharply than the proposal did. We are not comparing different model classes; rather, we are comparing EM and Bayesian inference for the same Gaussian HMM. On this dataset, Bayesian inference does not clearly improve clustering performance, but it does provide a principled posterior over transition persistence and emission parameters. That posterior is valuable when we want to understand whether long-duration states, such as sitting or laying, are estimated confidently or whether the apparent persistence is itself uncertain.

The proposal feedback also asked why uncertainty quantification should matter when the original motivation was label scarcity. The answer is that in a fully unsupervised sequential problem, uncertainty is one of the few ways to assess whether a learned latent explanation is stable. Even if the clustering metrics are similar, the Bayesian model exposes how much trust we should place in the transition dynamics.

At the same time, the results show the practical tradeoff clearly: Bayesian inference is more computationally expensive and only modestly different in predictive behavior from EM on this benchmark. A stronger extension would be to use the raw inertial signals directly, incorporate semi-supervised labels when available, or move to a nonparametric extension such as the HDP-HMM.
